\section{The curvature coefficient on homogeneous sources}
\label{curvatureobs:section}

The exact positive kernel improves the analytic estimate, but does not
make its numerical coefficient small for every source of a prescribed
pair of dimensions. We compute that coefficient on homogeneous dyadic
sources and exhibit an obstruction to every sequence of common selected
levels. The obstruction concerns the sufficient coefficient \(W_{n,m}\),
not the actual difference of affine densities or the actual distance law.

\subsection{An exact tail-sum formula}

Let \(K\subset[0,1]^2\) carry a homogeneous dyadic probability \(\mu\):
every occupied square of depth \(k\) has mass \(2^{-F(k)}\), and there
are \(2^{F(k)}\) such squares. Assume the same descendant counts in
every occupied square, as in the product-digit constructions above.
Thus an occupied depth-\(m\) square has exactly
\(2^{F(k)-F(m)}\) occupied depth-\(k\) descendants for \(k\ge m\).
Suppose
\[
 0\le F(k+1)-F(k)\le2.
\]
Use the unique admissible half-open coding; the constructions below have
infinitely many forced-zero levels, so this coding also applies to every
point of \(K\). Put
\begin{equation}\label{curvatureobs:tail}
 G(k)=F(k)-k,\qquad
 H(L)=\sum_{k\ge L}2^{-G(k)}.
\end{equation}
The following comparisons permit \(+\infty\).

\begin{lemma}[Homogeneous curvature coefficient]
\label{curvatureobs:formula}
For integers \(0\le n<m\), the coefficient of
Proposition~\ref{band:energy-criterion} satisfies
\begin{equation}\label{curvatureobs:exact}
 W^\mu_{n,m}\asymp
 2^{G(m)}H(\max\{m,2n\}).
\end{equation}
The comparison constants are absolute.
Moreover, for every probability \(\eta\) supported on \(K\),
\begin{equation}\label{curvatureobs:allsource}
 W^\eta_{n,m}\ge c\,2^{G(m)}H(\max\{m,2n\}).
\end{equation}
This lower bound also holds if the depth-\(m\) dyadic partition in
\(W^\eta_{n,m}\) is replaced by any countable measurable partition
whose pieces have diameter at most \(C_0 2^{-m}\); its constant may
then depend on \(C_0\). In particular
\begin{equation}\label{curvatureobs:longjump}
 m\ge2n\quad\Longrightarrow\quad W^\eta_{n,m}\ge c.
\end{equation}
\end{lemma}
\begin{proof}
Fix an occupied depth-\(m\) square \(Q\). For \(k\ge m\), the
conditional natural probability has pair-ball mass
\[
 (\mu_Q\times\mu_Q)\{|x-x'|\le2^{-k}\}
 \asymp 2^{F(m)-F(k)}.
\]
For the upper bound, a radius-\(2^{-k}\) ball meets only boundedly
many depth-\(k\) squares, each of the displayed conditional mass.
For the lower bound, pairs in the same depth-\((k+1)\) descendant
have distance at most \(\sqrt2\,2^{-k-1}<2^{-k}\); their probability
is \(2^{F(m)-F(k+1)}\), which is at least one quarter of the
displayed quantity.

The positive radial layer-cake formula for
\(k_\delta(u)=\min\{u^{-1},\delta^2u^{-3}\}\), with
\(\delta=2^{-2n}\), consequently gives
\begin{equation}\label{curvatureobs:dyadickernel}
 J_{2^{-2n}}(\mu_Q)
 \asymp
 2^{F(m)}
 \sum_{k\ge m}2^{-F(k)}
              \min\{2^k,2^{3k-4n}\}.
\end{equation}
The part of the layer cake at radii larger than \(2^{-m}\) is
bounded by a constant times the \(k=m\) summand and is harmless.
This follows either by integrating the kernel derivative there, or
by using pair mass at most one. The dyadic comparison of the
remaining integral uses the pair-ball estimates above.

When \(m<2n\), the terms \(2^{3k-4n-F(k)}\), \(m\le k<2n\),
increase by at least a factor two at every step. Their sum is bounded
by a constant times \(2^{2n-F(2n)}\), which is already the first
term of the remaining tail. When \(m\ge2n\), there is no such
initial part. Thus \eqref{curvatureobs:dyadickernel} reduces to
\[
 J_{2^{-2n}}(\mu_Q)
 \asymp 2^{F(m)}H(\max\{m,2n\}).
\]
These bounds are uniform in \(Q\). Multiplying their square roots
by the masses \(\mu(Q)\), which sum to one, proves
\eqref{curvatureobs:exact}.

Here is a lower bound independent of the chosen source probability.
Write \(L=\max\{m,2n\}\). For every \(x,x'\in K\),
\begin{equation}\label{curvatureobs:tree-kernel}
 k_{2^{-2n}}(x-x')
 \ge c\sum_{k\ge L}2^k
       \mathbf1_{\{x,x'\text{ lie in the same depth-}k
                              \text{ original square}\}}.
\end{equation}
If the right side is nonzero, then
\(|x-x'|\le\sqrt2\,2^{-L}\le\sqrt2\,2^{-2n}\);
on this region the left kernel is at least a constant times
\(|x-x'|^{-1}\). The geometric sum on the right is at most
\(C|x-x'|^{-1}\). On the diagonal both sides are infinite.
This proves \eqref{curvatureobs:tree-kernel}.

Let \(P\) be a positive-\(\eta\)-mass piece of a partition of diameter
at most \(C_0 2^{-m}\). It meets at most \(C(C_0)\) original
depth-\(m\) squares, hence at most
\[
 C(C_0)2^{F(k)-F(m)}
\]
original depth-\(k\) squares. For its normalized restriction \(\eta_P\),
Cauchy--Schwarz on these descendant masses gives a same-square pair
probability at least
\(c(C_0)2^{F(m)-F(k)}\). Integrating
\eqref{curvatureobs:tree-kernel} and using Tonelli yields
\[
 J_{2^{-2n}}(\eta_P)
 \ge c(C_0)2^{F(m)}H(L).
\]
Taking square roots, multiplying by \(\eta(P)\), and summing all pieces
proves \eqref{curvatureobs:allsource}. If \(m\ge2n\), the term \(k=m\)
in \(H(m)\) proves \eqref{curvatureobs:longjump}.
\end{proof}

\subsection{The limiting scalar quantity}

Suppose \(d=1+a>1\), \(F(k)\ge dk-C\), and scales \(N_j\to\infty\)
realize a continuous normalized profile:
\[
 \sup_{0\le x\le1}
 \left|\frac{F(N_jx)-N_jx}{N_j}-g(x)\right|\longrightarrow0,
 \qquad g(1)=a.
\]
Here \(F\) is interpolated linearly. For fixed
\(0<x<y\le1\), with \(2x\le1\), put
\(n_j=\lfloor xN_j\rfloor\), \(m_j=\lfloor yN_j\rfloor\).
Then
\begin{equation}\label{curvatureobs:profile-limit}
 \lim_{j\to\infty}\frac{\log_2 W^\mu_{n_j,m_j}}{N_j}
 =
 g(y)-\min_{\max\{y,2x\}\le z\le1}g(z).
\end{equation}
Indeed the terms of \(H\) up to \(N_j\) lie between their largest
term and \(N_j+1\) times their largest term. Uniform profile convergence
identifies its exponential rate. The remaining infinite tail is at most
\(C_a2^{-aN_j}\), by the lower Frostman barrier. Since \(g(1)=a\),
this cannot increase the limiting exponential rate. An integer closest
to a minimizing point supplies the matching lower bound. Equation
\eqref{curvatureobs:exact} completes the proof.

Thus a long jump cannot remove a future low value of \(F(k)-k\).
The next lemma makes this observation uniform over all selected-level
sequences and all source probabilities.

\subsection{A barrier for every common selected-level sequence}

\begin{lemma}[A crossing barrier]\label{curvatureobs:barrier}
Suppose there are integers \(L_j\to\infty\), integers \(u_j\ge2L_j\),
and a constant \(C_1<\infty\), independent of \(j\), such that
\begin{equation}\label{curvatureobs:barrier-hyp}
 \min_{L_j<k\le2L_j}G(k)\ge G(u_j)-C_1.
\end{equation}
For every probability \(\eta\) supported on \(K\), and every strictly
increasing sequence of common depths \(n_i\to\infty\), infinitely many
of its edges satisfy
\[
 W^\eta_{n_i,n_{i+1}}\ge c\,2^{-C_1}.
\]
The same conclusion holds for the bounded-diameter common-depth
partitions in Lemma~\ref{curvatureobs:formula}.
\end{lemma}
\begin{proof}
For large \(j\), choose the edge \(n_i\le L_j<n_{i+1}=m\).
If \(m\ge2n_i\), apply \eqref{curvatureobs:longjump}.
Otherwise
\[
 L_j<m<2n_i\le2L_j\le u_j.
\]
In \eqref{curvatureobs:allsource}, keep just the \(k=u_j\) term.
The result and \eqref{curvatureobs:barrier-hyp} give
\[
 W^\eta_{n_i,m}
 \ge c\,2^{G(m)-G(u_j)}\ge c\,2^{-C_1}.
\]
Every fixed edge crosses only finitely many \(L_j\), so infinitely
many distinct edges have this bound.
\end{proof}

In particular, every positive-power series in these coefficients
diverges. More generally any sufficient criterion requiring their
convergence to zero fails. This conclusion permits changing the source
measure entirely; it does not assume that it is dominated by the
original homogeneous measure.

\subsection{Examples with finite positive critical measures}

\begin{theorem}\label{curvatureobs:criticalsets}
For every
\[
 1<d<\frac32,\qquad 2d-1\le D<2,
\]
there is a compact homogeneous dyadic set \(K\) such that
\[
 0<\mathcal H^d(K)<\infty,\qquad
 0<\mathcal P^D(K)<\infty.
\]
In particular \(\HH K=d\) and \(\PP K=D\). Its natural probability
\(\mu\) is \(d\)-Frostman and equals both normalized critical measures:
\begin{equation}\label{curvatureobs:criticalmeasures}
 \mu=\frac{\mathcal H^d|_K}{\mathcal H^d(K)}
     =\frac{\mathcal P^D|_K}{\mathcal P^D(K)}.
\end{equation}
For every probability \(\eta\) supported on \(K\), every sequence
of increasing common levels, and every sequence of partitions at
those levels with the uniform diameter bound in
Lemma~\ref{curvatureobs:formula}, the resulting curvature coefficients
fail to tend to zero.
\end{theorem}
\begin{proof}
Put \(V(n)=2\lfloor Dn/2\rfloor\). Its increments are zero or two.
Choose a sufficiently large initial lower stopping depth, with an
even branching count within two of \(dn\) and strictly below \(V(n)\).
A finite initial path with increments zero or two realizes this count.

From each lower stop increase \(F\) by two at every step until first
meeting \(V\). This is finite because \(D<2\). The meeting is exact:
the even difference \(F-V\) starts negative and increases by zero
or two. After the meeting depth \(c\), follow \(F(n)=V(n)\) until
an arbitrarily large integer \(b\). Then keep \(F\) constant until
the next lower stop
\[
 a=\left\lceil V(b)/d\right\rceil.
\]
Choose \(b\) large enough that \(c<a/2<b<a\), and repeat. This
is possible since \(a/b\to D/d\in(1,2)\).
Apart from the finite initial path, the construction gives
\begin{equation}\label{curvatureobs:criticalbarriers}
 dn-2\le F(n)\le Dn,\qquad
 F(b)=Db+O(1),\qquad da-d<F(a)\le da.
\end{equation}
Use four children at an increment of two and only the lower-left
child at an increment of zero. The infinite flat phases give
unambiguous half-open addresses as above.

On \(a/2<k\le b\) the count is \(V(k)\), so
\[
 G(k)\ge(D-1)k-2
       \ge\frac{D-1}{2}a-2
       \ge(d-1)a-2\ge G(a)-2.
\]
On \(b\le k\le a\), the count is constant and
\(G(k)=F(a)-k\ge G(a)\). Taking \(L=\lfloor a/2\rfloor\)
and \(u=a\) gives \eqref{curvatureobs:barrier-hyp} with \(C_1=2\).
The lower stops tend to infinity. The crossing lemma proves the
claimed coefficient obstruction, including when \(D=2d-1\).

We verify the assertions about critical measures. The lower bound
in \eqref{curvatureobs:criticalbarriers}, and bounded overlap of
balls with comparable dyadic cells, imply the \(d\)-Frostman bound.
Consequently \(\mathcal H^d(K)>0\). The complete cylinder covers
at the lower stops have bounded \(d\)-dimensional cost, proving
\(\mathcal H^d(K)<\infty\).

The upper bound \(F(n)\le Dn+C\) gives
\[
 \mu(B(x,r))\ge c r^D\qquad(x\in K,\ 0<r\le1).
\]
Indeed a cylinder containing \(x\) at a sufficiently fine comparable
depth lies wholly inside \(B(x,r)\), and has mass at least
\(c r^D\). Any disjoint ball packing centered on \(K\) therefore
has bounded sum of \(D\)-th powers of diameters. This proves finite
\(D\)-dimensional packing premeasure and hence
\(\mathcal P^D(K)<\infty\).

At each peak \(b\), the natural measure of every ball of radius
\(C2^{-b}\), for fixed \(C\), is at most \(C'2^{-Db}\).
For any subset \(A\subset K\), take a maximal separated family in
\(A\) at this scale. Its enlarged balls cover \(A\), so its
cardinality is at least \(c\mu^*(A)2^{Db}\). The smaller disjoint
balls have total \(D\)-dimensional cost at least \(c\mu^*(A)\).
Letting the peak depths tend to infinity shows that the packing
premeasure of every such \(A\) is at least \(c\mu^*(A)\).
Applying this to the intersections with \(K\) of every countable
cover of \(K\), and using outer subadditivity, proves
\(\mathcal P^D(K)>0\).

At any fixed depth the cylinders are disjoint compact Borel subsets
of \(K\), and are translations of one common scaled tail set.
Translation invariance makes the restriction of each finite positive
critical measure give them equal masses. After normalization these
masses are exactly \(2^{-F(n)}\). The cylinders generate the Borel
sigma algebra of \(K\); uniqueness of finite measures proves
\eqref{curvatureobs:criticalmeasures}. Finally \(d>1\) gives
\(I_1(\mu)<\infty\), so the obstruction is not caused by failure of
that prerequisite.
\end{proof}

\subsection{Compatibility with the sharp profile examples}

The obstruction can hold even inside the proved positive-distance region.
For example, take \(d=11/10\) and \(D=121/100<28/23=B_{\rm H}(d)\)
in Theorem~\ref{curvatureobs:criticalsets}.
Theorem~\ref{target:hardgap} gives a positive-length self-pinned distance
set on this support, whereas its curvature coefficients fail for every
source and every sequence of common depths. Thus failure of this
coefficient test cannot itself be a necessary obstruction to that conclusion.

For \(0<a<b\le1\), the single-collapse profile
\[
 g(x)=\min\{bx,1+a-x\}
\]
has a lower-barrier valley at \(v=1\): \(g(v)=av\).
The two-collapse profile of Section~\ref{hardgap:two-witness} has
\[
 g(x)=\min\{bx,av+|x-v|,1+a-x\},\qquad v=t\in(0,1).
\]
It too satisfies \(g(v)=av\). In either case
\begin{equation}\label{curvatureobs:valley}
 g(x)>av\qquad(av/b<x<v).
\end{equation}
For the two-collapse profile, each of the three entries in the minimum
is strictly above \(av\) on this interval. For the single-collapse
profile the same check uses its two entries.

Suppose first that \(b>2a\). Choose
\[
 \frac{av}{b}<c<\frac v2.
\]
By \eqref{curvatureobs:valley}, the compact interval \([c,2c]\)
has minimum strictly greater than \(av\).
Use the fixed-measure realization from
Theorem~\ref{realization:theorem}, and set
\[
 L_j=\lfloor cN_j\rfloor,\qquad u_j=\lceil vN_j\rceil.
\]
Uniform convergence of the realized profiles gives, for some
\(\epsilon>0\) and all sufficiently large \(j\),
\[
 \min_{L_j<k\le2L_j}G(k)-G(u_j)\ge\epsilon N_j.
\]
Thus Lemma~\ref{curvatureobs:barrier} applies, uniformly over every
replacement source probability on this one fixed support.

At the critical equality \(b=2a<1\), uniform convergence alone need
not give a constant error. The exact catch-up realization supplies it.
Write \(h(x)=x+g(x)\), \(d=1+a\), and \(D=1+b<2\).
At a preceding endpoint \(M\), assume a continuous branching function
has value \(\mathcal F(M)=dM\). Choose the next endpoint \(N\) much
larger than \(M\) and put on this block
\begin{equation}\label{curvatureobs:catchup}
 \mathcal F(t)=
 \min\{Nh(t/N),\,2t-(2-d)M\},\qquad M\le t\le N.
\end{equation}
The initial piece of \(h\) is \(Dx\). The entries meet there at
\[
 t_0=\frac{2-d}{2-D}M.
\]
Choose \(N\) so large that this meeting lies in the initial linear
piece and \(t_0<vN/2\). Before \(t_0\) the second entry is smaller;
afterward it remains at least the first because the template is
2-Lipschitz. Hence \(\mathcal F(t)=Nh(t/N)\) on
\([vN/2,N]\). This block joins continuously, has slopes in \([0,2]\),
satisfies \(dt\le\mathcal F(t)\le Dt\), and ends at \(dN\).

Choose successive \(M/N\to0\) and discretize by
\(F(k)=2\lfloor\mathcal F(k)/2\rfloor\). The increments are zero
or two and the rounding error is less than two. The same argument as
in Theorem~\ref{realization:theorem} gives a fixed \(d\)-Frostman
homogeneous source with exact dimensions \(d,D\), realizing the chosen
profile. At
\[
 L_j=\lfloor vN_j/2\rfloor,\qquad u_j=\lceil vN_j\rceil,
\]
the template has \(g(x)\ge av\) on \([v/2,v]\), since \(b=2a\).
The bounded rounding and endpoint errors therefore prove
\eqref{curvatureobs:barrier-hyp} with a constant independent of \(j\).
The barrier applies at critical equality as well.

For the single-collapse matching range the limiting profile cost is
\(C_1(a,b)\), and for the two-collapse matching range it is \(C_2(a,b)\),
as proved in Section~\ref{hardgap:section}. Whenever the relevant cost
is at least \(a\) and \(b\ge2a\), the same fixed support therefore has
both features: its limiting profile cost is at least \(a\), so no
fixed strict cost margin holds along the realizing scales, and
every source probability on it fails the curvature-coefficient
convergence test for every sequence of common levels. The continuity
of the limiting profile cost was proved in
Section~\ref{hardgap:compactness}, so the realizations preserve that
assertion.

This is a limitation of these separate numerical tests. It does not
rule out cancellation between comparison increments, a new geometric
estimate, or a new argument mixing the two approximation schemes.
In particular it does not assert that the actual laws have a divergent
variation, that selected source--pin pairs cannot help, or that the
distance sets have zero length. Arbitrary source-dependent trees with
different comparison scales on different branches are also outside
the common-level assertion proved here.
